\subsection{Sampling and statistical inference}

The entries in this section use one invented sample: twelve monthly active returns of a manager against a benchmark, in percent, $1.2, -0.4, 0.8, 2.1, -1.5, 0.6, 1.9, -0.2, 0.3, 1.4, -0.9, 1.1$, with mean 0.5333, sample standard deviation 1.1106 and standard error $1.1106/\sqrt{12}=0.3206$. The comparison of variances uses a second, independent sample: fifteen values drawn by computer, independently of the first sample, from a normal distribution with mean 0 and standard deviation 0.5 and rounded to one decimal, $-0.2$, 0.1, 0.3, 0.4, 0.8, 0.2, 0.4, 0.4, 0.3, 0.1, $-0.7$, 0.0, $-0.3$, 1.5, 0.3, with sample standard deviation 0.4968. They illustrate the mechanics and are not evidence about any manager.

\begin{qfmethod}[Sampling, the central limit theorem and confidence intervals: estimand, sampling distribution, exact $t$ and chi-square intervals, normal approximations and the repeated-sampling meaning]\label{meth:confidence-intervals}
\qffield{Statement} If $X_1,\dots,X_n$ are independent and identically distributed with finite mean $\mu$ and variance $0<\sigma^2<\infty$, the standardised mean $(\bar X-\mu)/(\sigma/\sqrt n)$ converges in distribution to a standard normal (the central limit theorem); for finite $n$ the normal is an approximation whose quality depends on the distribution, with $n\ge30$ a common rule of thumb \citep[sec.~7.1.2]{probability-statistics-random-processes}. For a normal sample with unknown variance, $T=(\bar X-\mu)/(S/\sqrt n)$ has a $t$ distribution with $n-1$ degrees of freedom and $(n-1)S^2/\sigma^2$ a chi-square distribution with $n-1$ degrees of freedom; they give the intervals $\bar X\pm t_{\alpha/2,n-1}S/\sqrt n$ for $\mu$ and $\bigl[(n-1)S^2/\chi^2_{\alpha/2,n-1},(n-1)S^2/\chi^2_{1-\alpha/2,n-1}\bigr]$ for $\sigma^2$, valid for any $n$ (sec.~8.3.3). A 95 percent interval means that ``in the long run about 95\,\% of these intervals would contain the true mean''; it ``does not mean that there is a 95\,\% probability that the interval contains the true mean'' \citep[sec.~1.3.5.2]{statistical-methods-handbook}.

Samples can be drawn in several ways: a \emph{simple random sample} gives every sample of the same size ``an equal chance of being selected''; a \emph{stratified} sample divides the population into strata and takes ``a proportionate number from each stratum''; a \emph{cluster} sample randomly selects some groups; a \emph{systematic} sample takes every $n$th item from a random start; \emph{convenience} sampling uses ``results that are readily available'' and invites sampling bias, where ``some members of the population are not as likely to be chosen as others''. Sampling itself causes sampling error \citep[sec.~1.2]{introductory-statistics-2023}. In finance, a time series of returns is not a random sample from a population of periods, and a sample of surviving funds is a convenience sample (Interpretation).

\qffield{Notation} \emph{Estimand}: the population quantity sought, here the manager's mean monthly active return $\mu$. \emph{Sampling design}: how the data arose; the formulas above assume a random sample, independent draws from one distribution. $\bar X$ the estimator; $S/\sqrt n$ its estimated standard error; $\chi^2_{q,k}$ the upper-$q$ point in the source's notation. Interpretation: \emph{exact} means the stated coverage holds for every $n$ under the model (the $t$ and chi-square intervals for normal i.i.d.\ data); \emph{approximate} means it holds only as $n$ grows (the normal interval $\bar X\pm1.96S/\sqrt n$ for non-normal data).

\qffield{Conditions} Monthly returns are rarely independent or identically distributed: volatility clusters and means drift. With positive serial correlation $S/\sqrt n$ understates the standard error of the mean, and the interval is too narrow; time-series standard errors (\hyperref[meth:sharpe-ratio-standard-error]{Sharpe-ratio standard errors}) or a block bootstrap (\hyperref[meth:bootstrap-resampling]{bootstrap}) are then needed. The chi-square interval for a variance depends on normality much more than the $t$ interval for a mean does (Interpretation). The interval is about the estimand of the design: a sample of one manager's months says nothing about other managers.

\qffield{Practice} Example. The exact 95 percent $t$ interval for $\mu$ is $0.5333\pm2.2010\times0.3206$, from $-0.17$ to 1.24 percent a month; the normal-approximation interval, $\pm1.96$ standard errors, is $-0.10$ to 1.16 and too narrow for twelve observations. The 95 percent interval for the variance is 0.62 to 3.56, a standard deviation of 0.79 to 1.89 percent. In 40{,}000 simulated normal samples of twelve, 95.1 percent of the $t$ intervals covered the true mean. In AR(1) samples with coefficient 0.5 the same formula covered only 75 percent.

\qffield{Pitfalls} Reading a computed interval as having a 95 percent probability of containing $\mu$. Drawing conclusions about a population of funds or stocks from a sample selected by survival or availability. Applying i.i.d.\ formulas to overlapping or autocorrelated returns. Treating a narrow interval from many highly correlated observations as precise.

\qffield{See also} \hyperref[meth:hypothesis-testing]{hypothesis tests}; \hyperref[def:return-sample-statistics]{sample statistics}; \hyperref[meth:skill-luck-significance-tests]{skill against luck}.

\end{qfmethod}

\begin{qfmethod}[Hypothesis tests: null and alternative, test statistic, significance level, $p$-value, Type I and Type II errors, power; $t$, chi-square and $F$ tests]\label{meth:hypothesis-testing}
\qffield{Statement} A test sets a null $H_0$ (for example $\mu=\mu_0$, or $\mu\le\mu_0$) against an alternative $H_1$ ($\mu\neq\mu_0$, or $\mu>\mu_0$). A Type I error rejects $H_0$ when it is true, and the test is designed so that its probability is at most $\alpha$; a Type II error accepts $H_0$ when $H_1$ is true, with probability $\beta(\mu)$ depending on the true value. For a mean, $W=(\bar X-\mu_0)/(S/\sqrt n)$ is compared with $t_{\alpha/2,n-1}$ for normal data, an exact test, or with $z_{\alpha/2}$ for non-normal data with large $n$, an approximation; a two-sided test accepts exactly when $\mu_0$ lies in the corresponding confidence interval \citep[sec.~8.4.3]{probability-statistics-random-processes}. The $p$-value is ``the lowest significance level $\alpha$ that results in rejecting the null hypothesis'' (sec.~8.4.4). A variance is tested with $T=(N-1)(s/\sigma_0)^2$ against chi-square critical values with $N-1$ degrees of freedom, and two variances with $F=s_1^2/s_2^2$ against $F$ critical values with $N_1-1$ and $N_2-1$ degrees of freedom \citep[secs.~1.3.5.8--1.3.5.9]{statistical-methods-handbook}, under the conditions that ``each population is normally distributed'' and the samples are ``independent samples from the 2 populations'' \citep[Lesson~12, sec.~12.1, ``Test to Compare Two Population Variances'']{applied-statistics-course-notes}; the Levene test ``is less sensitive than the Bartlett test to departures from normality'' (sec.~1.3.5.10).

\qffield{Notation} Power $=1-\beta$, the probability of rejecting $H_0$ when a specified alternative is true. Interpretation: for the one-sided $t$ test at level $\alpha$, with a true mean $\delta\sigma$ above $\mu_0$, the statistic has a noncentral $t$ distribution with noncentrality $\delta\sqrt n$, so power is $P(T_{n-1,\delta\sqrt n}>t_{\alpha,n-1})$, exact for normal data.

\qffield{Conditions} The stated $\alpha$ is the Type I error rate only if the test's assumptions hold and the hypothesis was fixed before seeing the data; testing many variants and reporting the best inflates it (\hyperref[meth:multiple-testing-control]{multiple-testing control}). Failing to reject is not evidence that $H_0$ is true, particularly when power is low. The chi-square and $F$ tests for variances assume normal data and are sensitive to departures from it (Interpretation, consistent with the Levene comparison above). The $F$ test also needs the two samples to be independent of each other; normality within each sample does not supply that. Returns of two managers or strategies over the same months share market shocks and are usually correlated, and the ordinary $F$ reference distribution then does not apply. Interpretation, by simulation: with equal true variances and a correlation of 0.9 between paired months ($n=12$), 40{,}000 simulated two-sided 5 percent tests reject 0.05 percent of the time; with a true ratio of standard deviations of 1.5 they reject 6 percent of the time, against 24 percent for independent samples. Paired series need a test that models the pairing, which this entry does not cover. Statistical significance is not economic size.

\qffield{Practice} Example. $H_0:\mu\le0$ against $H_1:\mu>0$ at 5 percent: $t=0.5333/0.3206=1.664$ with 11 degrees of freedom, a one-sided $p$-value of 0.062 (two-sided 0.124), so $H_0$ is not rejected, and 0 lies inside the 95 percent interval. If the true mean were half a standard deviation, the test would detect it with probability only 0.49 at $n=12$, 0.90 at $n=36$ and 0.99 at $n=60$. Under a true null, 4.9 percent of 40{,}000 simulated tests rejected. For the first sample against the independent second sample, $F=1.1106^2/0.4968^2=5.00$ on 11 and 14 degrees of freedom, beyond the upper two-sided 5 percent critical value of 3.09 ($p=0.006$), so equal variances are rejected. The arithmetic can be checked against software, but only the sampling design justifies the $F$ distribution. Testing $\sigma=1$ for the first gives $T=13.57$ on 11 degrees of freedom, $p=0.52$.

\qffield{Pitfalls} Reading $p$ as the probability that $H_0$ is true. Treating ``not significant'' as ``no skill'' with a low-power test (\hyperref[meth:skill-luck-significance-tests]{skill against luck} gives the data needs). Using an $F$ test on fat-tailed returns without checking normality, or on two return series from the same months as if they were independent.

\qffield{See also} \hyperref[meth:confidence-intervals]{confidence intervals}; \hyperref[meth:nonparametric-tests]{nonparametric tests}; \hyperref[def:financial-distributions]{distributions}.

\end{qfmethod}

\begin{qfmethod}[Choosing between parametric and nonparametric tests: the sign test and the signed-rank test against the $t$ test, their assumptions and power]\label{meth:nonparametric-tests}
\qffield{Statement} The signed-rank test is a one-sample or paired test that replaces the $t$ test's normality assumption by symmetry. Its assumptions are that each difference has a symmetric distribution, that the differences are mutually independent and share the same mean, and that they are measured on at least an interval scale. It ``is not as powerful as the $t$-test when the distributional assumptions are in fact valid''. It is more powerful than the sign test because it uses magnitudes as well as signs, but it assumes more, so the sign test is preferred when symmetry is doubtful. Its statistic is the smaller rank sum, with tabulated critical values for $N\le30$ and, for larger $N$, a normal approximation with mean $n(n+1)/4$ and standard deviation $\sqrt{n(n+1)(2n+1)/24}$ \citep[``Signed Rank Test'']{signed-rank-test-reference-2023}.

\qffield{Notation} Interpretation: the \emph{sign test} counts positive observations; under $H_0$ that the median is zero, with independent observations, the count is binomial with $p=1/2$, so its $p$-value is an exact binomial tail. The signed-rank statistic $W^+$ sums the ranks of $|x_i|$ over positive $x_i$; its exact null distribution comes from the $2^n$ equally likely sign patterns.

\qffield{Conditions} Nonparametric does not mean assumption-free: both tests need independent observations, the sign test tests a median, and the signed-rank test needs symmetry for its null to be about the centre. Interpretation: choose by the question and the data. Use the $t$ test for a mean when normality is plausible or $n$ is large; use the signed-rank test when tails are heavy but roughly symmetric; use the sign test when only direction is trusted or symmetry fails. None repairs serial dependence.

\qffield{Practice} Example, the twelve active returns. Eight are positive: the exact one-sided sign-test $p$-value is $\sum_{k\ge8}\binom{12}k/2^{12}=0.194$. The ranks of the positive values sum to $W^+=58$; among all 4{,}096 sign patterns, 7.6 percent give 58 or more, the exact one-sided $p$-value; the normal approximation with continuity correction gives 0.073. The $t$ test gave 0.062. Each test uses less information than the one before it, and none rejects at 5 percent.

\qffield{Pitfalls} Calling a rank test ``distribution-free'' and dropping its independence and symmetry assumptions. Switching to whichever test gives the smallest $p$-value after seeing the results. Testing a mean hypothesis with a median test.

\qffield{See also} \hyperref[meth:hypothesis-testing]{hypothesis tests}; \hyperref[meth:bootstrap-resampling]{bootstrap} (another route when formulas are doubtful).

\end{qfmethod}

\begin{qfmethod}[Testing correlation and independence: the $t$ test of zero correlation, Spearman's rank correlation with ties, the chi-square test of independence, and overlapping observations]\label{meth:correlation-independence-tests}
\qffield{Statement} \emph{Pearson correlation.} To test $H_0:\rho=0$ against $\rho\ne0$, use $t=r\sqrt{n-2}/\sqrt{1-r^2}$ with $n-2$ degrees of freedom; $t$ has the sign of $r$, and the $p$-value is the area in both tails. Equivalently, $r$ is significant at 5 percent when it lies beyond $\pm t_c/\sqrt{t_c^2+n-2}$, which is 0.602 for $n=11$ and 0.456 for $n=19$ \citep[sec.~12.4]{introductory-statistics-2023}. The test assumes a linear relation, normal and equally dispersed $y$ at each $x$, independent residuals and a random sample.

\emph{Spearman's rank correlation} replaces each observation by its rank, giving tied values their average rank, and computes the Pearson correlation of the ranks. It measures monotone, not only linear, association. NIST tests $\rho_s=0$ with tabulated quantiles for $n\le30$ without ties, and otherwise with the normal approximation $\rho_s\sqrt{n-1}$ \citep[``Rank Correlation Independence Test'']{nist-dataplot-rank-correlation}.

\emph{Independence in a contingency table.} With observed counts $O$ and expected counts $E=(\text{row total})(\text{column total})/(\text{grand total})$, the statistic $\sum(O-E)^2/E$ has a chi-square distribution with $(r-1)(c-1)$ degrees of freedom under independence. The test is right-tailed and needs every expected count to be at least five (OpenStax, sec.~11.3).

\qffield{Notation} $r$ sample correlation; $n$ pairs; $\rho_s$ Spearman's coefficient; $r\times c$ the table's rows and columns. Interpretation: Spearman's coefficient is unchanged by a strictly increasing transform of either variable. A strictly decreasing transform of one variable changes only its sign, and of both leaves it unchanged. The shortcut $1-6\sum d_i^2/(n(n^2-1))$, with $d_i$ the rank differences, equals it only without ties. A weakly monotone transform can create ties and change the value.

\qffield{Conditions} The $t$ test needs independent pairs. In quant research the \emph{information coefficient} of a signal, its correlation with the next period's return, is tested this way, and the \emph{rank IC} is the Spearman version, robust to outliers. When forward returns overlap, as with a daily signal and 20-day returns, consecutive pairs share most of their returns. If the signal is persistent, the effective sample is far smaller than $n$, and the nominal test rejects far too often (Interpretation, checked below). Use non-overlapping periods or standard errors that allow for the overlap.

\qffield{Practice} \emph{Reproduced}: OpenStax's exam data, $r=0.6631$ with $n=11$, give $t=2.658$ and $p=0.026$; the critical values for $n=6$ to 19 are reproduced. In Example~11.6, 839 volunteers classified by type and hours (a $3\times3$ table) give $\chi^2=12.9909$ with 4 degrees of freedom and $p=0.0113$, so hours and type are dependent at 5 percent (sec.~11.3).

\emph{Checked} (synthetic):
\begin{itemize}
\item On 20{,}000 independent normal samples of 30 pairs, the 5 percent Pearson test rejected 4.84 percent of the time. Spearman's test with the $t$ approximation on midranks rejected 4.95 percent, and with NIST's normal approximation 4.78 percent.
\item A tied fixture of 12 pairs has $\rho_s=0.9288$ on midranks, while the shortcut gives 0.9301.
\item A persistent signal (AR(1), coefficient 0.95) correlated with independent returns over 500 days: the 5 percent test rejected 6.5 percent of the time with one-day returns and 62.9 percent with overlapping 20-day returns.
\end{itemize}

\qffield{Pitfalls} Reading a significant $r$ as a large or causal effect. Using the $d_i^2$ shortcut with ties. Running a chi-square test with expected counts below five. Testing an IC built from overlapping returns with $n-2$ degrees of freedom.

\qffield{See also} \hyperref[meth:hypothesis-testing]{hypothesis testing}; \hyperref[meth:nonparametric-tests]{nonparametric tests}; \hyperref[def:return-dependence]{dependence}.

\end{qfmethod}
